% GENERATED FILE. Edit canonical lessons, then run npm run generate:books.
% canonical-source-hash: fnv1a64:e37bb1bcf6d00661
% canonical-lessons: ML-C474-nastamakuka, ML-C474-reyilppalam, ML-C474-catiyirannuka, ML-C474-tunninilkkuka, ML-C474-payuka

\chapter{Missing the Train}
\label{ch:ml-missing-the-train}
\hlchaptermodality{474}

\begin{chapteropening}
\textbf{By the end of this chapter:} \emph{I can say to miss, a railway track, to jump off, to hang on and to speed along.}

Complete the last lesson of chapter 474: I can say to miss, a railway track, to jump off, to hang on and to speed along.
\end{chapteropening}

\section[\texorpdfstring{\ml{നഷ്ടമാകുക} (naṣṭamākuka)}{നഷ്ടമാകുക (naṣṭamākuka)}]{\ml{നഷ്ടമാകുക} (naṣṭamākuka) — to miss; to be lost}

\label{lesson:ML-C474-nastamakuka}



\begin{quote}
Before the new one: say the Malayalam for to shove, then the Malayalam for to rush up.
\end{quote}

\subsection*{You'll want to know: \ml{നഷ്ടമാകുക}}
\textbf{\ml{നഷ്ടമാകുക}} — \emph{naṣṭamākuka} — \textquotedblleft{}to miss; to be lost\textquotedblright{}.

\begin{quote}
\ml{അഞ്ച്} \ml{മിനിറ്റ്} \ml{വൈകിയതുകൊണ്ട്} \ml{എനിക്ക്} \ml{ബസ്} \ml{നഷ്ടമായി}.
\end{quote}

\begin{quote}
\emph{añcŭ miniṟṟŭ vaikiyatukoṇṭŭ enikkŭ basŭ naṣṭamāyi.}
\end{quote}

\begin{quote}
Because I was five minutes late, I missed the bus.
\end{quote}

Past: \ml{നഷ്ടമായി}, \emph{naṣṭamāyi}.

The one who misses takes -\ml{ക്ക്}: \ml{എനിക്ക്} \ml{തീവണ്ടി} \ml{നഷ്ടമായി}, I missed the train. A chance can be lost the same way.

\begin{cousinweb}
\textbf{\ml{നഷ്ടമാകുക}} (\emph{naṣṭamākuka}) is a Malayalam verb built from \ml{നഷ്ടം}, a loss, borrowed from Sanskrit \emph{naṣṭa}, \textquotedblleft{}lost, perished\textquotedblright{} (CDIAL 7027), and \ml{ആകുക}, to become.
\end{cousinweb}

\begin{grammarlens}[title={What you've built}]
One more word for getting around.
\end{grammarlens}

\subsection*{Your turn}
\begin{itemize}
\raggedright
  \item \emph{Say it:} \emph{naṣṭamākuka}
  \item \emph{Say it:} \emph{naṣṭamākuka}, once more
  \item \emph{Say it:} \emph{añcŭ miniṟṟŭ vaikiyatukoṇṭŭ enikkŭ basŭ naṣṭamāyi.} — Because I was five minutes late, I missed the bus.
  \item \emph{Recall:} say the Malayalam for to shove, then the Malayalam for to rush up, then say \emph{naṣṭamākuka} again
\end{itemize}

\begin{tcolorbox}[breakable,colback=teal!4,colframe=teal!35!black,title={Before you move on}]
What does \ml{നഷ്ടമാകുക} mean? (\textquotedblleft{}To miss; to be lost\textquotedblright{}.) Say it once more.
\end{tcolorbox}
\section[\texorpdfstring{\ml{റെയിൽപ്പാളം} (ṟeyilppāḷaṁ)}{റെയിൽപ്പാളം (ṟeyilppāḷaṁ)}]{\ml{റെയിൽപ്പാളം} (ṟeyilppāḷaṁ) — a railway track}

\label{lesson:ML-C474-reyilppalam}



\begin{quote}
Before the new one: say the Malayalam for to rush up, then the Malayalam for to miss.
\end{quote}

\subsection*{You'll want to know: \ml{റെയിൽപ്പാളം}}
\textbf{\ml{റെയിൽപ്പാളം}} — \emph{ṟeyilppāḷaṁ} — \textquotedblleft{}a railway track\textquotedblright{}.

\begin{quote}
\ml{റെയിൽപ്പാളം} \ml{മുറിച്ചുകടക്കുന്നത്} \ml{അപകടകരമാണ്}.
\end{quote}

\begin{quote}
\emph{ṟeyilppāḷaṁ muṟiccukaṭakkunnatŭ apakaṭakaramāṇŭ.}
\end{quote}

\begin{quote}
Crossing the railway track is dangerous.
\end{quote}

In speech \ml{പാളം} alone often means the whole railway track.

\begin{cousinweb}
\textbf{\ml{റെയിൽപ്പാളം}} (\emph{ṟeyilppāḷaṁ}) is a Malayalam compound built from \ml{റെയിൽ}, which was borrowed from English \emph{rail}, and the native word \ml{പാളം}, a cast metal bar (DEDR 4114).
\end{cousinweb}

\begin{grammarlens}[title={What you've built}]
One more word for getting around.
\end{grammarlens}

\subsection*{Your turn}
\begin{itemize}
\raggedright
  \item \emph{Say it:} \emph{ṟeyilppāḷaṁ}
  \item \emph{Say it:} \emph{ṟeyilppāḷaṁ}, once more
  \item \emph{Say it:} \emph{ṟeyilppāḷaṁ muṟiccukaṭakkunnatŭ apakaṭakaramāṇŭ.} — Crossing the railway track is dangerous.
  \item \emph{Recall:} say the Malayalam for to rush up, then the Malayalam for to miss, then say \emph{ṟeyilppāḷaṁ} again
\end{itemize}

\begin{tcolorbox}[breakable,colback=teal!4,colframe=teal!35!black,title={Before you move on}]
What does \ml{റെയിൽപ്പാളം} mean? (\textquotedblleft{}A railway track\textquotedblright{}.) Say it once more.
\end{tcolorbox}
\section[\texorpdfstring{\ml{ചാടിയിറങ്ങുക} (cāṭiyiṟaṅṅuka)}{ചാടിയിറങ്ങുക (cāṭiyiṟaṅṅuka)}]{\ml{ചാടിയിറങ്ങുക} (cāṭiyiṟaṅṅuka) — to jump off}

\label{lesson:ML-C474-catiyirannuka}



\begin{quote}
Before the new one: say the Malayalam for to miss, then the Malayalam for a railway track.
\end{quote}

\subsection*{You'll want to know: \ml{ചാടിയിറങ്ങുക}}
\textbf{\ml{ചാടിയിറങ്ങുക}} — \emph{cāṭiyiṟaṅṅuka} — \textquotedblleft{}to jump off\textquotedblright{}.

\begin{quote}
\ml{ബസ്} \ml{നിൽക്കുന്നതിന്} \ml{മുൻപ്} \ml{ഒരു} \ml{കുട്ടി} \ml{ചാടിയിറങ്ങി}.
\end{quote}

\begin{quote}
\emph{basŭ nilkkunnatinŭ munpŭ oru kuṭṭi cāṭiyiṟaṅṅi.}
\end{quote}

\begin{quote}
A child jumped off before the bus stopped.
\end{quote}

Past: \ml{ചാടിയിറങ്ങി}, \emph{cāṭiyiṟaṅṅi}.

\ml{ചാടി}, having jumped, plus \ml{ഇറങ്ങുക}, to get down.

\begin{cousinweb}
\textbf{\ml{ചാടിയിറങ്ങുക}} (\emph{cāṭiyiṟaṅṅuka}) is a native Dravidian compound of \ml{ചാടുക}, to leap (DEDR 2440), and \ml{ഇറങ്ങുക}, to get down (DEDR 516).
\end{cousinweb}

\begin{grammarlens}[title={What you've built}]
One more word for getting around.
\end{grammarlens}

\subsection*{Your turn}
\begin{itemize}
\raggedright
  \item \emph{Say it:} \emph{cāṭiyiṟaṅṅuka}
  \item \emph{Say it:} \emph{cāṭiyiṟaṅṅuka}, once more
  \item \emph{Say it:} \emph{basŭ nilkkunnatinŭ munpŭ oru kuṭṭi cāṭiyiṟaṅṅi.} — A child jumped off before the bus stopped.
  \item \emph{Recall:} say the Malayalam for to miss, then the Malayalam for a railway track, then say \emph{cāṭiyiṟaṅṅuka} again
\end{itemize}

\begin{tcolorbox}[breakable,colback=teal!4,colframe=teal!35!black,title={Before you move on}]
What does \ml{ചാടിയിറങ്ങുക} mean? (\textquotedblleft{}To jump off\textquotedblright{}.) Say it once more.
\end{tcolorbox}
\section[\texorpdfstring{\ml{തൂങ്ങിനിൽക്കുക} (tūṅṅinilkkuka)}{തൂങ്ങിനിൽക്കുക (tūṅṅinilkkuka)}]{\ml{തൂങ്ങിനിൽക്കുക} (tūṅṅinilkkuka) — to hang on; to cling}

\label{lesson:ML-C474-tunninilkkuka}



\begin{quote}
Before the new one: say the Malayalam for a railway track, then the Malayalam for to jump off.
\end{quote}

\subsection*{You'll want to know: \ml{തൂങ്ങിനിൽക്കുക}}
\textbf{\ml{തൂങ്ങിനിൽക്കുക}} — \emph{tūṅṅinilkkuka} — \textquotedblleft{}to hang on; to cling\textquotedblright{}.

\begin{quote}
\ml{തിരക്കുള്ള} \ml{ബസിൽ} \ml{ചില} \ml{യാത്രക്കാർ} \ml{വാതിലിൽ} \ml{തൂങ്ങിനിന്നു}.
\end{quote}

\begin{quote}
\emph{tirakkuḷḷa basil cila yātrakkār vātilil tūṅṅininnu.}
\end{quote}

\begin{quote}
On the crowded bus some passengers hung on at the door.
\end{quote}

Past: \ml{തൂങ്ങിനിന്നു}, \emph{tūṅṅininnu}.

\ml{തൂങ്ങി}, hanging, plus \ml{നിൽക്കുക}, to stay so: to keep hanging on.

\begin{cousinweb}
\textbf{\ml{തൂങ്ങിനിൽക്കുക}} (\emph{tūṅṅinilkkuka}) is a native Dravidian compound of \ml{തൂങ്ങുക}, to hang (DEDR 3376), and \ml{നിൽക്കുക}, to stand (DEDR 3675).
\end{cousinweb}

\begin{grammarlens}[title={What you've built}]
One more word for getting around.
\end{grammarlens}

\subsection*{Your turn}
\begin{itemize}
\raggedright
  \item \emph{Say it:} \emph{tūṅṅinilkkuka}
  \item \emph{Say it:} \emph{tūṅṅinilkkuka}, once more
  \item \emph{Say it:} \emph{tirakkuḷḷa basil cila yātrakkār vātilil tūṅṅininnu.} — On the crowded bus some passengers hung on at the door.
  \item \emph{Recall:} say the Malayalam for a railway track, then the Malayalam for to jump off, then say \emph{tūṅṅinilkkuka} again
\end{itemize}

\begin{tcolorbox}[breakable,colback=teal!4,colframe=teal!35!black,title={Before you move on}]
What does \ml{തൂങ്ങിനിൽക്കുക} mean? (\textquotedblleft{}To hang on; to cling\textquotedblright{}.) Say it once more.
\end{tcolorbox}
\section[\texorpdfstring{\ml{പായുക} (pāyuka)}{പായുക (pāyuka)}]{\ml{പായുക} (pāyuka) — to speed along; to rush}

\label{lesson:ML-C474-payuka}



\begin{quote}
Before the new one: say the Malayalam for to jump off, then the Malayalam for to hang on.
\end{quote}

\subsection*{You'll want to know: \ml{പായുക}}
\textbf{\ml{പായുക}} — \emph{pāyuka} — \textquotedblleft{}to speed along; to rush\textquotedblright{}.

\begin{quote}
\ml{ആംബുലൻസ്} \ml{ആശുപത്രിയിലേക്ക്} \ml{പാഞ്ഞു}.
\end{quote}

\begin{quote}
\emph{āmbulansŭ āśupatriyilēkkŭ pāññu.}
\end{quote}

\begin{quote}
The ambulance sped to the hospital.
\end{quote}

Past: \ml{പാഞ്ഞു}, \emph{pāññu}.

The past is \ml{പാഞ്ഞു}: a verb in -\ml{ആയുക} makes its past in -\ml{ആഞ്ഞു}.

\begin{cousinweb}
\textbf{\ml{പായുക}} (\emph{pāyuka}) is native Dravidian, \textquotedblleft{}to leap, to rush out\textquotedblright{} (DEDR 4087).
\end{cousinweb}

\begin{grammarlens}[title={What you've built}]
One more word for getting around.
\end{grammarlens}

\subsection*{Your turn}
\begin{itemize}
\raggedright
  \item \emph{Say it:} \emph{pāyuka}
  \item \emph{Say it:} \emph{pāyuka}, once more
  \item \emph{Say it:} \emph{āmbulansŭ āśupatriyilēkkŭ pāññu.} — The ambulance sped to the hospital.
  \item \emph{Recall:} say the Malayalam for to jump off, then the Malayalam for to hang on, then say \emph{pāyuka} again
\end{itemize}

\begin{tcolorbox}[breakable,colback=teal!4,colframe=teal!35!black,title={Before you move on}]
What does \ml{പായുക} mean? (\textquotedblleft{}To speed along; to rush\textquotedblright{}.) Say it once more.
\end{tcolorbox}
